% ============================================================================
% Appendix D — Example Pipeline Output
% ============================================================================
\chapter{Example Pipeline Output}
\label{app:examples}

This appendix presents annotated output from a complete \versimux{} evaluation run on a representative test page. All data structures are reproduced from actual pipeline JSON output.


% ---- D.1 Input ---------------------------------------------------------------
\section{Input HTML}
\label{app:input-html}

\begin{figure}[htbp]
  \centering
  \framebox[\textwidth]{\rule{0pt}{5cm}\parbox{0.9\textwidth}{\centering\large\sffamily [Figure D.1: Example Input HTML Page]\\ \vspace{0.8em} \normalsize\normalfont Screenshot of a sample login page used as input to the \versimux{} pipeline. The page includes an email input, a password input, a submit button, and a ``forgot password'' link, with several intentional accessibility violations: missing \code{<label>} elements, no \code{lang} attribute on the \code{<html>} tag, and insufficient colour contrast on the submit button.}}
  \caption{Example input page with planted accessibility violations.}
  \label{fig:example-input}
\end{figure}


% ---- D.2 Supervisor Output ---------------------------------------------------
\section{Supervisor Analysis Output}
\label{app:supervisor-output}

\begin{figure}[htbp]
  \centering
  \framebox[\textwidth]{\rule{0pt}{6cm}\parbox{0.9\textwidth}{\centering\large\sffamily [Figure D.2: Supervisor UIAnalysis JSON Output]\\ \vspace{0.8em} \normalsize\normalfont Formatted JSON output showing the \code{UIAnalysis} object: \code{ui\_purpose} = ``User login'', \code{ui\_type} = ``login form'', \code{accessibility\_risk\_level} = ``high'', \code{detected\_issues\_hint} listing 5 specific element-level violations (missing labels, missing \code{lang}, no error feedback region, missing \code{autocomplete}), \code{interactive\_elements} listing 4 elements with their selectors and \code{is\_accessible} flags, and \code{critical\_paths} defining the login flow.}}
  \caption{Example supervisor \code{UIAnalysis} output.}
  \label{fig:example-uianalysis}
\end{figure}


% ---- D.3 Persona Traces -----------------------------------------------------
\section{Persona Interaction Traces}
\label{app:persona-traces}

\begin{figure}[htbp]
  \centering
  \framebox[\textwidth]{\rule{0pt}{7cm}\parbox{0.9\textwidth}{\centering\large\sffamily [Figure D.3: Example Persona Action Trace]\\ \vspace{0.8em} \normalsize\normalfont Complete 8-step action trace from a ``Screen Reader User'' persona. Shows the sequence of actions: (1)~observe page structure, (2)~type email into \code{\#email} input, (3)~type password into \code{\#password} input, (4)~click \code{.btn-login}, (5)~observe error feedback (none visible), (6)~observe --- issue detected: missing form error announcement, (7)~attempt to navigate back, (8)~stop (max steps). Each step includes the \code{action\_type}, \code{target\_selector}, \code{reasoning}, \code{success} flag, and any triggered \code{issue\_id}.}}
  \caption{Example persona interaction trace (Screen Reader User).}
  \label{fig:example-trace}
\end{figure}


% ---- D.4 Clustering Output ---------------------------------------------------
\section{Issue Clusters}
\label{app:clustering-output}

\begin{figure}[htbp]
  \centering
  \framebox[\textwidth]{\rule{0pt}{5cm}\parbox{0.9\textwidth}{\centering\large\sffamily [Figure D.4: Example Clustering Output]\\ \vspace{0.8em} \normalsize\normalfont JSON showing 3 clusters generated from 9 verified issues: \code{cluster\_1} (``Missing form field labels'', 4 issues, dominant severity: high), \code{cluster\_2} (``No error feedback mechanism'', 3 issues, dominant severity: critical), \code{cluster\_3} (``Colour contrast violations'', 2 issues, dominant severity: medium). Each cluster includes \code{affected\_personas}, \code{affected\_elements}, and \code{representative\_description}.}}
  \caption{Example issue clustering output.}
  \label{fig:example-clusters}
\end{figure}


% ---- D.5 Patches and Conflicts -----------------------------------------------
\section{Patch Proposals and Conflict Resolution}
\label{app:patches-output}

\begin{figure}[htbp]
  \centering
  \framebox[\textwidth]{\rule{0pt}{7cm}\parbox{0.9\textwidth}{\centering\large\sffamily [Figure D.5: Example Patch Proposals and Conflict Resolution]\\ \vspace{0.8em} \normalsize\normalfont Shows: (1)~3 \code{PatchProposal} objects generated by recommender agents (one per cluster), including \code{patch\_type}, \code{target\_element}, \code{before\_snippet}/\code{after\_snippet} for the HTML label fix, \code{css\_snippet} for the colour contrast fix, and \code{js\_snippet} for the error feedback fix. (2)~One \code{ConflictRecord} where two patches targeted overlapping selectors, and the resulting \code{NegotiationSession} with the mediator's resolution.}}
  \caption{Example patch proposals and conflict resolution output.}
  \label{fig:example-patches}
\end{figure}


% ---- D.6 Verification Result ------------------------------------------------
\section{Verification Result}
\label{app:verification-output}

\begin{figure}[htbp]
  \centering
  \framebox[\textwidth]{\rule{0pt}{5cm}\parbox{0.9\textwidth}{\centering\large\sffamily [Figure D.6: Example Verification Result]\\ \vspace{0.8em} \normalsize\normalfont Per-persona \code{VerificationResult} objects showing: for ``Screen Reader User'' --- 4 issues before patching, 3 resolved, 1 remaining (low severity), 0 regressions, \code{fully\_resolved} = true. For ``Keyboard-Only User'' --- 3 issues before, 3 resolved, 0 remaining, \code{fully\_resolved} = true. Final verification decision: \code{verification\_passed} = true.}}
  \caption{Example per-persona verification results.}
  \label{fig:example-verification}
\end{figure}


% ---- D.7 Final Report -------------------------------------------------------
\section{Generated Report}
\label{app:report-output}

\begin{figure}[htbp]
  \centering
  \framebox[\textwidth]{\rule{0pt}{6cm}\parbox{0.9\textwidth}{\centering\large\sffamily [Figure D.7: Example Diagnostic Report Summary]\\ \vspace{0.8em} \normalsize\normalfont Formatted excerpt from the final \code{DiagnosticReport}: \code{overall\_score} = 8.5, \code{total\_issues\_found} = 9, \code{issues\_resolved\_count} = 7, \code{issues\_remaining\_count} = 2, \code{regressions\_introduced} = 0, \code{verification\_passed} = true. The \code{executive\_summary} and \code{top\_recommendations} fields are shown with their full narrative text.}}
  \caption{Example final diagnostic report summary.}
  \label{fig:example-report}
\end{figure}
